\section*{Cómo leer este documento}
\addcontentsline{toc}{section}{Cómo leer este documento}
\markboth{Cómo leer este documento}{Cómo leer este documento}

Este documento tiene cuatro partes. Se pueden leer por separado.

\begin{description}
  \item[Parte I. El informe.] Es el informe final del equipo, tal como está en el repositorio
    (\cod{docs/INFORME_FINAL.md}): qué es Alicanto, qué quedó hecho, cómo se midió, el resultado, sus
    límites y cómo trabajamos. Son unas doce páginas y no suponen conocimientos previos.
  \item[Parte II. Cómo funciona por dentro.] Explica cada pieza con el detalle del código: qué recibe,
    qué calcula, con qué fórmula y con qué parámetros. Es para quien quiera revisar el algoritmo o
    reimplementarlo.
  \item[Parte III. Cómo se midió.] Describe los datos, el protocolo de los tres equipos, las métricas y
    las pruebas estadísticas, y da los resultados completos, incluido el motivo por el que cada familia
    de ataques se frena o se escapa.
  \item[Parte IV. Apéndices.] Tablas de parámetros, el modelo ajustado, formatos de archivo, el mapa
    del repositorio y los comandos para repetir cada cifra.
\end{description}

\subsection*{De dónde salen las cifras}

\begin{itemize}
  \item \textbf{Las tablas y figuras con datos no se escribieron a mano.} Las genera el script
    \cod{docs/informe/hacer_datos.py} a partir de los archivos versionados (\cod{results/},
    \cod{models/}, \cod{redteam/} y las corridas grabadas de la demo). Con \cod{--check}, el script
    vuelve a generarlas y avisa si dejaron de coincidir con el repositorio.
  \item \textbf{Parte del detalle es nuevo.} El repositorio guarda los totales de la medición con las
    marcas del proxy, pero no el detalle llamada a llamada. Para este informe se volvió a correr el
    detector congelado sobre el set de prueba, por el mismo camino de código, guardando además la
    evidencia de cada llamada y la regla de cada rotación. El script comprueba que esa corrida
    reproduce los totales publicados. No se ajustó nada con ella: el set ya estaba abierto y el
    detector, congelado.
  \item \textbf{Cada cifra del texto dice de qué archivo sale,} y el \apref{ap:repetir} trae el comando
    que la repite.
\end{itemize}

\subsection*{Palabras que usamos con un sentido preciso}

\begin{description}
  \item[Frenado.] Un ataque está frenado (en el código, \cod{prevented}) si el atacante no llegó a
    sacar más de 96 de los 160 bits de una misma versión de la clave. No significa que no haya salido
    nada: significa que lo que salió no alcanza.
  \item[Marcas.] Cada llamada lleva dos marcas: si saca datos hacia afuera (\cod{is_sink}) y si su
    resultado es contenido no confiable (\cod{untrusted}). Los documentos del repositorio las llaman
    «banderas». Importan porque el resultado principal cambia según quién las calcula
    (\cref{sec:hallazgo-marcas}).
  \item[Sumidero.] Llamada que saca datos hacia afuera: una petición HTTP, un issue o un comentario
    en GitHub, o un \cod{git push}.
  \item[Equipos rojo, azul y morado.] Son tres papeles con reglas escritas, no tres grupos de
    personas. El rojo fue una sesión de agente aparte, con la regla de no abrir el detector.
\end{description}

\begin{limite}
\textbf{Dos cifras que no se deben confundir.} Con las marcas que traían las grabaciones, Alicanto
frenó 180 de 180 ataques. Con las marcas que calcula el proxy en vivo, que es como funcionaría
instalado, frenó 120 de 180. La cifra que decimos es la segunda. La primera solo aparece en este
documento junto a esta explicación.
\end{limite}

Los números van con coma decimal y punto de miles. Las horas son de Santiago.
\clearpage
